\subsection{Femtosecond Laser Micromachining}
While bonding focuses on joining substrates, vapor cells and 
many applications require the precise removal 
of glass to create cavities, channels, or complex 3D structures. 
Due to the transparency, brittleness, and chemical resistance of 
glass, traditional machining often fails to meet high-precision 
requirements. \textbf{CITA}
\begin{itemize}
    \item \textbf{Mechanical Machining:} Utilizing diamond-tipped 
    tools or ultrasonic drilling. While effective for macro-scale 
    features, it introduces \textit{micro-cracks} and significant 
    internal stress, which can lead to structural failure.
    \item \textbf{Wet Chemical Etching:} Typically involves 
    Hydrofluoric Acid (HF). This process is \textit{isotropic}, 
    meaning it etches in all directions at once. Consequently, 
    it is difficult to create deep, narrow channels with high 
    aspect ratios.
    \item \textbf{Reactive Ion Etching (RIE):} A dry etching 
    process using plasma. It offers high precision but is limited 
    by extremely slow etch rates (micrometers per hour), making 
    it suitable only for surface patterns or thin-film glass.
\end{itemize}
Femtosecond Laser Writing (FLW) has emerged as a key
technology because it allows for "sub-surface" machining 
without affecting the glass exterior. 
\begin{itemize}
    \item \textbf{Nonlinear Absorption:} Glass is transparent 
    to most light; however, a femtosecond laser delivers energy 
    in ultrashort pulses ($10^{-15}$ s). At the focal point, 
    the intensity is high enough to trigger \textit{multiphoton absorption},
    allowing the laser to "write" inside the glass without 
    damaging the surface.
    \item \textbf{FLICE Process:} Femtosecond Laser Irradiation 
    followed by Chemical Etching (FLICE) is a two-step technique:
    \begin{enumerate}
        \item \textbf{Irradiation:} The laser modifies the 
        local molecular structure of the glass, increasing 
        its chemical reactivity.
        \item \textbf{Etching:} The glass is submerged in 
        an etchant (e.g., KOH or HF), which dissolves the 
        laser-modified tracks up to 1000 times faster than 
        the surrounding bulk glass.
    \end{enumerate}
    \item \textbf{Zero Heat-Affected Zone (HAZ):} Because 
    the pulse duration is shorter than the thermal diffusion 
    time of the material, the process is considered 
    "cold machining." This prevents melting, burrs, or thermal cracking.
\end{itemize}
Femtosecond laser technology is uniquely suited for the optics and semiconductor 
industries because it overcomes 
the isotropic limitations of chemical etching and the mechanical stresses 
of traditional drilling. It remains the 
only viable method for creating "buried" 3D micro-channels and complex 
internal cavities within transparent substrates.
\textbf{MILLORAR. POSAR PAPER DE GIANVITO (FOTOS I EXPLICACIONS)}